\documentclass[12pt,a4paper]{ctexart}
\usepackage{geometry}\geometry{margin=2.2cm}
\usepackage{fancyhdr,hyperref,longtable,booktabs,enumitem,listings,xcolor}
\lstset{basicstyle=\ttfamily\small,breaklines=true,frame=single}
\pagestyle{fancy}\fancyhf{}\fancyhead[L]{OPC DA 自主采集操作手册}\fancyhead[R]{V1.0}\fancyfoot[C]{\thepage}
\title{\textbf{OPC DA 自主采集操作手册}\\LONGNAME $\to$ TagID $\to$ ReadList 全链路}
\author{约翰·刘}\date{2026年7月13日}

\begin{document}\maketitle
\tableofcontents
\newpage

\section{概述}

本手册覆盖油田作业区 IO 服务器 292 台设备、46,046 个测点的完整 OPC DA 数据采集方案。

\subsection{适用范围}
\begin{itemize}[leftmargin=*]
  \item 目标: pSpace 6.0 实时数据库 (11.66.12.130:8889)
  \item 环境: 32 位 Python 3.11 + psAPISDK.dll
  \item 数据源: Oracle PROJECT\_TAGPAR $\to$ parse.db tag\_definition
\end{itemize}

\subsection{采集链路}
\begin{verbatim}
Step 1: parse.db 查 LONGNAME
  SELECT longname FROM tag_definition
  WHERE driver_name='OPC_FC_Client' AND device_id='DX1ZRZ'
  → /CY1C8K/DX1ZRZ/ZD010833ALA172

Step 2: psAPI 解析 TagID
  psAPI_Tag_GetIdByLongName(handle, LONGNAME, &tagId)
  → TagID = 11945

Step 3: 读取实时值
  psAPI_Real_ReadList(handle, 1, [11945], &values, &errors)
  → value = 2.2017, timestamp = 2026-07-13 22:01:10, quality = GOOD
\end{verbatim}

\section{环境准备}

\subsection{32 位 Python}
\begin{lstlisting}
C:\Python311-32\python.exe --version
# Python 3.11.9 (32-bit)
\end{lstlisting}

\subsection{依赖文件}
\begin{itemize}[leftmargin=*]
  \item psAPISDK.dll + 所有依赖 DLL (IO ServerOnLine 目录)
  \item parse.db (本地 tag\_definition 表，46,046 行)
\end{itemize}

\subsection{网络要求}
\begin{itemize}[leftmargin=*]
  \item 可达 11.66.12.130:8889 (pSpace)
  \item Oracle 可通过 WinRM 131 中转 (可选，仅需导入时)
\end{itemize}

\section{设备清单查询}

\subsection{按驱动查询设备}
\begin{lstlisting}[language=Python]
import sqlite3
db = sqlite3.connect('data/parse.db')

# 列出所有驱动
for row in db.execute("""
    SELECT driver_name, COUNT(DISTINCT device_id) as devs, COUNT(*) as tags
    FROM tag_definition GROUP BY driver_name
"""):
    print(f"{row[0]}: {row[1]} devices, {row[2]} tags")

# 列出某驱动的设备
for row in db.execute("""
    SELECT DISTINCT device_id, well_name
    FROM tag_definition WHERE driver_name='OPC_FC_Client'
    LIMIT 10
"""):
    print(f"  {row[0]} -> {row[1]}")
\end{lstlisting}

\subsection{按设备查询 LONGNAME}
\begin{lstlisting}[language=Python]
device = 'DX1ZRZ'
for row in db.execute("""
    SELECT tag_name, longname, data_type, range_max, range_min, unit
    FROM tag_definition
    WHERE device_id = ?
    LIMIT 10
""", (device,)):
    print(f"{row[0]:30s} {row[1]:50s} [{row[4]}-{row[3]}] {row[5]}")
\end{lstlisting}

\section{LONGNAME $\to$ TagID 解析}

\subsection{标准工具}
\begin{lstlisting}
C:\Python311-32\python.exe tools/pspace_collector.py --ids "11945,5657"
\end{lstlisting}

\subsection{批量解析}
\begin{lstlisting}[language=Python]
import ctypes, os
os.add_dll_directory(r'D:\ai\dgiot_lite\io服务器分析\IO ServerOnLine')
dll = ctypes.CDLL('psAPISDK.dll')

def resolve_tags(paths):
    dll.psAPI_Common_StartAPI.restype = ctypes.c_int32
    dll.psAPI_Common_StartAPI()
    dll.psAPI_Server_Connect.argtypes = [ctypes.c_char_p]*3 + [ctypes.POINTER(ctypes.c_uint16)]
    dll.psAPI_Server_Connect.restype = ctypes.c_int32
    h = ctypes.c_uint16(0xFFFF)
    dll.psAPI_Server_Connect(b'11.66.12.130:8889', b'admin', b'DQYTA11_pass', ctypes.byref(h))

    dll.psAPI_Tag_GetIdByLongName.argtypes = [ctypes.c_uint16, ctypes.c_char_p, ctypes.POINTER(ctypes.c_uint32)]
    dll.psAPI_Tag_GetIdByLongName.restype = ctypes.c_int32

    results = {}
    for path in paths:
        tid = ctypes.c_uint32(0)
        r = dll.psAPI_Tag_GetIdByLongName(h, path, ctypes.byref(tid))
        if r == 0:
            results[path.decode()] = tid.value

    dll.psAPI_Server_Disconnect(h)
    dll.psAPI_Common_StopAPI()
    return results
\end{lstlisting}

\section{实时数据读取}

\subsection{单次读取}
\begin{lstlisting}[language=Python]
def read_tags(tag_ids):
    dll = ctypes.CDLL('psAPISDK.dll')
    dll.psAPI_Common_StartAPI.restype = ctypes.c_int32
    dll.psAPI_Common_StartAPI()
    dll.psAPI_Server_Connect.argtypes = [ctypes.c_char_p]*3 + [ctypes.POINTER(ctypes.c_uint16)]
    dll.psAPI_Server_Connect.restype = ctypes.c_int32
    h = ctypes.c_uint16(0xFFFF)
    dll.psAPI_Server_Connect(b'11.66.12.130:8889', b'admin', b'DQYTA11_pass', ctypes.byref(h))

    dll.psAPI_Real_ReadList.argtypes = [ctypes.c_uint16, ctypes.c_uint32, ctypes.c_void_p, ctypes.c_void_p, ctypes.c_void_p]
    dll.psAPI_Real_ReadList.restype = ctypes.c_int32

    n = len(tag_ids)
    ids_arr = (ctypes.c_uint32 * n)(*tag_ids)
    ppV = ctypes.c_void_p(0)
    r = dll.psAPI_Real_ReadList(h, n, ids_arr, ctypes.byref(ppV), 0)

    results = []
    if r == 0 and ppV.value:
        raw = ctypes.string_at(ppV.value, n * 24)
        for i, tid in enumerate(tag_ids):
            import struct
            v = struct.unpack('<f', raw[i*24+16:i*24+20])[0]
            ts = struct.unpack('<I', raw[i*24:i*24+4])[0]
            quality = raw[i*24+20]
            results.append({'tag_id': tid, 'value': v, 'timestamp': ts, 'quality': quality})

    dll.psAPI_Server_Disconnect(h)
    dll.psAPI_Common_StopAPI()
    return results
\end{lstlisting}

\subsection{PS\_DATA 内存布局}
\begin{table}[H]\centering
\begin{tabular}{lll}
\toprule
\textbf{偏移} & \textbf{大小} & \textbf{字段} \\
\midrule
0-3 & 4B & Time.Seconds (uint32 LE) \\
4-7 & 4B & Time.NanoSec \\
8-11 & 4B & DataType \\
12-15 & 4B & (reserved) \\
16-19 & 4B & \textbf{Float Value} \\
20-23 & 4B & Quality (0xC0=GOOD) \\
\bottomrule
\end{tabular}
\end{table}

\section{常用 SQL 查询}

\begin{table}[H]\centering
\begin{tabular}{lp{8cm}}
\toprule
\textbf{场景} & \textbf{SQL} \\
\midrule
查设备所有 Tag & \texttt{SELECT * FROM tag\_definition WHERE device\_id='DX1ZRZ'} \\
按驱动统计 & \texttt{SELECT driver\_name, COUNT(*) FROM tag\_definition GROUP BY 1} \\
查某 Tag 的 LONGNAME & \texttt{SELECT longname FROM tag\_definition WHERE tag\_name LIKE '\%ALA172'} \\
查量程范围 & \texttt{SELECT device\_id, tag\_name, range\_min, range\_max, unit FROM tag\_definition WHERE driver\_name='OPC\_FC\_Client' LIMIT 10} \\
\bottomrule
\end{tabular}
\end{table}

\section{数据覆盖}

\begin{table}[H]\centering
\begin{tabular}{lrrrl}
\toprule
\textbf{驱动} & \textbf{设备} & \textbf{Tag} & \textbf{访问} & \textbf{状态} \\
\midrule
OPC\_FC\_Client (Kepware) & 34 & 40,956 & LONGNAME$\to$TagID$\to$ReadList & 已验证 \\
IM\_A11\_RTU (油田) & 161 & 3,638 & 同上 & 已验证 \\
Standard\_Umodbus & 50 & 1,452 & 同上 & 已验证 \\
保护装置 & 47 & -- & ReadList 直读 & 已验证 \\
\bottomrule
\end{tabular}
\end{table}

\section{限制与注意事项}

\begin{itemize}[leftmargin=*]
  \item \textbf{单次连接限制}: 每次 Connect 只能执行一次 ReadList (最多约 20 个 Tag)
  \item \textbf{解析和读取分开}: Tag\_GetIdByLongName 和 Real\_ReadList 必须在不同连接中
  \item \textbf{限流}: 连续连接间隔 $>$ 5 秒，避免 -19998
  \item \textbf{32 位必须}: psAPISDK.dll 仅支持 32 位 Python
  \item \textbf{网络依赖}: 需可达 11.66.12.130:8889
  \item \textbf{离线可用}: parse.db 本地查询无需 Oracle
\end{itemize}

\section{故障排查}

\begin{table}[H]\centering
\begin{tabular}{llp{5cm}}
\toprule
\textbf{错误码} & \textbf{含义} & \textbf{解决} \\
\midrule
-19998 & 未初始化 & 重新 Connect 获取新 handle \\
-19997 & Tag 数超限 & 减少至 20 个以内 \\
-18591 & 密码错误 & 确认凭据 \\
-18592 & 用户不存在 & 确认用户名 \\
-19391 & Tag 不存在 & 检查 LONGNAME 是否正确 \\
\bottomrule
\end{tabular}
\end{table}

\end{document}
